\subsection{双曲線関数の逆関数}

$(\sinh u)'=\cosh u>0$なので,\,$\sinh u$は実数全体で狭義単調増加する.また$u\to\infty$で$\sinh u\to\infty$,\,$u\to-\infty$で$\sinh u\to-\infty$だから,どの実数$x$も$x=\sinh u$と一意に書ける.この逆関数を$\operatorname{arsinh}x$と書く.

$x=\sinh u$なら,\,$\cosh u>0$と$\cosh^2u=1+\sinh^2u$より$\cosh u=\sqrt{1+x^2}$である.定義式を足すと$\sinh u+\cosh u=\e^u$だから
\begin{equation}\label{式:逆双曲線正弦}
u=\operatorname{arsinh}x=\log\bigl(x+\sqrt{1+x^2}\bigr)
\end{equation}
となる.対数の中は,\,$\sqrt{1+x^2}>|x|$によって常に正である.逆関数の微分から
\[
\frac{\dd u}{\dd x}=\frac1{\cosh u}
=\frac1{\sqrt{1+x^2}}
\]
を得る.ここでも積分の値から逆関数を定義したのではなく,指数関数による定義から順に導いた.

この関係を置換積分に使おう.\,$a>0$として$x=a\sinh u$と置くと
\[
\sqrt{a^2+x^2}=a\cosh u,\qquad
\dd x=a\cosh u\dd u
\]
となり,根号と$\dd x$に同じ$\cosh u$が現れる.三角置換$x=a\tan\theta$では,根号が$a\sec\theta$,\,$\dd x$が$a\sec^2\theta\dd\theta$となって次数がずれたが,双曲線関数では次数がずれない.

\begin{bookpoint}
\textbf{POINT}\par
$\sqrt{a^2+x^2}$を含む積分は$x=a\sinh u$と置き,\,$\sqrt{a^2+x^2}=a\cosh u$,\,$\dd x=a\cosh u\dd u$を代入する.
\end{bookpoint}

\begin{example}{双曲線正弦への置換}
$\int_0^3\frac1{\sqrt{4+x^2}}\dd x$を求めよ.
\end{example}

\paragraphbf{方針}
$x=2\sinh u$と置くと,分母と$\dd x$の$2\cosh u$が打ち消し合い,\,$1$の積分になる.

\paragraphbf{解答}
$\log\frac{3+\sqrt{13}}2$

\paragraphbf{解法}
$x=2\sinh u$と置けば
\[
\dd x=2\cosh u\dd u,\qquad
\sqrt{4+x^2}=2\cosh u
\]
である.端点$x=0$は$u=0$へ,\,$x=3$は
\[
u=\operatorname{arsinh}\frac32
=\log\left(\frac32+\sqrt{1+\frac94}\right)
=\log\frac{3+\sqrt{13}}2
\]
へ移る.したがって
\[
\int_0^3\frac1{\sqrt{4+x^2}}\dd x
=\int_0^{\log\frac{3+\sqrt{13}}2}1\dd u
=\log\frac{3+\sqrt{13}}2
\]
となる.

\paragraphbf{三角置換との関係}
同じ積分を$x=2\tan\theta$で置換すると,\,$\int\sec\theta\dd\theta$が残る.\,\sectionref{節:相反三角関数}の結果を使えば
\[
\int\sec\theta\dd\theta=\log(\sec\theta+\tan\theta)+C
\]
である.ここで$\tan\theta=\frac x2$,\,$\sec\theta=\frac{\sqrt{4+x^2}}2$を代入すると,対数の中は$\frac{x+\sqrt{4+x^2}}2$となり,\,$\operatorname{arsinh}\frac x2$と一致する.2つの置換は$\tan\theta=\sinh u$で結ばれ,このとき$\sec\theta=\cosh u$,\,$\sec\theta+\tan\theta=\e^u$となる.正割の積分に現れた$\sec\theta+\tan\theta$は,双曲線関数で書けば$\e^u$だったのだ.

また\chapterref{章:部分積分}では,\,$\log\bigl(x+\sqrt{1+x^2}\bigr)$を微分すると$\frac1{\sqrt{1+x^2}}$になることを直接確かめた.式\eqref{式:逆双曲線正弦}は,その対数が$\operatorname{arsinh}x$であることを表している.

\paragraphbf{偶関数の逆関数}
$\cosh u$は偶関数なので,実数全体では逆関数を持たない.しかし$u\geqq0$に制限すれば,\,$1$から無限大へ狭義単調増加する.\,$x=\cosh u$を使う置換では,このように枝を選ぶ.\,$\sinh$と$\cosh$で,逆向きに戻すときの条件は同じではない.

\begin{example}{双曲線正接への置換}
$\int_0^{\frac12}\frac1{1-x^2}\dd x$を求めよ.
\end{example}

\paragraphbf{方針}
$x=\tanh u$と置くと,\,$\dd x=(1-x^2)\dd u$となって被積分関数の分母が消える.

\paragraphbf{解答}
$\frac12\log3$

\paragraphbf{解法}
$\tanh u=\frac{\e^{2u}-1}{\e^{2u}+1}$は実数全体で狭義単調増加し,値の範囲は$(-1,1)$である.\,$x=\tanh u$と置けば,式\eqref{式:双曲線正接の微分}によって
\[
\dd x=(1-\tanh^2u)\dd u=(1-x^2)\dd u
\]
となる.元へ戻す式は,\,$x=\frac{\e^{2u}-1}{\e^{2u}+1}$を$\e^{2u}$について解いて
\begin{align*}
x(\e^{2u}+1)&=\e^{2u}-1\\
\e^{2u}&=\frac{1+x}{1-x}\\
u&=\frac12\log\frac{1+x}{1-x}
\end{align*}
である.したがって端点$x=0,\frac12$は$u=0,\frac12\log3$へ移り,
\[
\int_0^{\frac12}\frac1{1-x^2}\dd x
=\int_0^{\frac12\log3}1\dd u
=\frac12\log3
\]
となる.

\paragraphbf{部分分数分解との関係}
部分分数分解で
\[
\frac1{1-x^2}=\frac12\left(\frac1{1+x}+\frac1{1-x}\right)
\]
と分けても解ける.原始関数は
\[
\frac12\log(1+x)-\frac12\log(1-x)
=\frac12\log\frac{1+x}{1-x}
\]
であり,\,$-1<x<1$で上の$u$と同じ式である.部分分数分解で現れる2つの対数の差が,双曲線正接の逆関数そのものになっている.
